\section{Automorphisms under elementary extension}

We prove that elementary extension preserves the entire weak homotopy
type of marked definable automorphism families.

Let $\mathcal M\preccurlyeq\mathcal N$ be an elementary extension of
o-minimal structures in the same language, both expanding real closed
fields.  Definability over either structure permits all of its
parameters.  For a closed bounded definable set $D\subset M^n$, let
$\mathcal S$ be a finite list of definable subsets of $D$, with each
subset designated either setwise preserved or pointwise fixed.  The
marked subsets need not be closed.  Write
\[
 G_{\mathcal M}(D;\mathcal S)
   =\operatorname{Aut}_{\mathrm{def},\mathcal M}(D;\mathcal S)
\]
for the simplicial group whose $q$-simplices are definable
homeomorphisms of $D\times\Delta^q_M$ preserving the entire parameter
projection and satisfying the designated conditions on
$S\times\Delta^q_M$ for every $S\in\mathcal S$.  Both the total map
and its inverse are jointly continuous.  Faces and degeneracies are
affine parameter pullbacks, and multiplication is composition.

\begin{theorem}\label{thm:elementary-extension-marked-automorphisms}
Extension of defining formulas induces a weak equivalence
\[
 G_{\mathcal M}(D;\mathcal S)
 \longrightarrow
 G_{\mathcal N}(D_{\mathcal N};\mathcal S_{\mathcal N}).
\]
This includes the empty set, for which both groups are trivial.
In particular, for any closed bounded definable manifold $B$ over
$\mathcal M$, it gives a weak equivalence
\[
 P_{\mathrm{def},\mathcal M}(B)
 \longrightarrow P_{\mathrm{def},\mathcal N}(B_{\mathcal N}),
\]
where the bottom of $B\times[0,1]$ is fixed pointwise and its free
top is preserved setwise.

If $B$ is the rational realization of a finite abstract closed PL
manifold, the square with the PL
coefficient-extension map commutes.  In each homotopy degree,
injectivity, surjectivity, or bijectivity of the PL-to-definable
comparison over $\mathcal M$ is equivalent to the corresponding
property over $\mathcal N$.  Thus the comparison is a weak
equivalence over one structure exactly when it is so over the other.
\end{theorem}

We first isolate the continuity assertion required to move
coefficients.  Elementary transfer alone would not provide this
movement.

\begin{lemma}\label{lem:coefficient-cells-joint-evaluation}
Fix a nonempty closed bounded definable set $E$ over $\mathcal M$
and a formula $\varphi(a;x,y)$ with fixed parameters in
$\mathcal M$.  Let $T$ be the definable set of tuples $a$ for which
the formula defines a homeomorphism $h_a:E\to E$ satisfying any
specified finite list of first-order conditions.  There is a finite
partition of $T$ into cells defined over $\mathcal M$ such that
both maps
\[
 (a,x)\longmapsto h_a(x),\qquad
 (a,x)\longmapsto h_a^{-1}(x)
\]
are jointly continuous on every cell times $E$.  The same partition
and continuity assertions hold after elementary extension.
\end{lemma}

\begin{proof}
Homeomorphism validity is first order: functionality and bijectivity
are first-order graph conditions, and continuity of the map and its
inverse are expressed by epsilon--delta formulas in all coordinates.
Thus the stated set $T$ is definable; quantifier elimination for the
o-minimal theory is not required.

Consider the graph family
\[
 \Gamma=\{(a,x,y):a\in T,\ x,y\in E,\ \varphi(a;x,y)\}.
\]
Definable Hardt triviality, followed by cell decomposition of the
base, gives a finite cell partition with a trivialization over each
cell.  The precise classical inputs are
\cite[Theorem~5.22, pp.~62--63; Theorem~2.10, p.~19]{Coste1999}.
These statements apply to a fixed o-minimal expansion of an arbitrary
real closed field; the graph family need not be proper over its base.
Choose $a_0\in C(M)$ for a nonempty cell $C$.  By identifying the
trivializing fibre with the fibre at $a_0$, write the trivialization as
\[
 \theta:C\times\Gamma_{a_0}\xrightarrow{\ \cong\ }\Gamma|_C
\]
over $C$.  The set $\Gamma_{a_0}$ is closed and bounded, since it is
the graph of a continuous map from the closed bounded set $E$ to itself.

Compose $\theta$ with the projection to the input coordinate.  The
result is a continuous definable bijection
\[
 J:C\times\Gamma_{a_0}\longrightarrow C\times E.
\]
We verify inverse continuity without using ordinary compactness of
a non-Archimedean field.  If $J^{-1}$ were discontinuous at
$z=(a,x)$, there would be $\epsilon>0$ such that $z$ lies in the
closure of the definable set of points $z'$ with
$\|J^{-1}(z')-J^{-1}(z)\|\geq\epsilon$.
Definable curve selection provides a curve $z(t)$ in this set tending
to $z$ as $t\downarrow0$.  Write
$J^{-1}(z(t))=(a(t),\gamma(t))$.  Its first coordinate tends to $a$,
and its second lies in the closed bounded set $\Gamma_{a_0}$.
Coordinatewise monotonicity gives a limit
$\gamma\in\Gamma_{a_0}$: bounded definable one-variable functions
have finite one-sided limits, and closedness retains the limit in
the set.  Continuity of $J$ gives $J(a,\gamma)=z$.  Bijectivity
then identifies $(a,\gamma)$ with $J^{-1}(z)$, contradicting the
chosen distance bound.  The curve-selection and definable-compactness
inputs are \cite[Theorems~3.2 and~3.4, pp.~26--27]{Coste1999}.

Consequently the output projection of $\theta J^{-1}$ is joint
continuous evaluation on $C\times E$.  Repeating the argument with
the input and output interchanged proves joint inverse evaluation.
All the objects just constructed and their continuity assertions are
defined over $\mathcal M$.  Elementarity transfers them to
$\mathcal N$.
\end{proof}

\begin{proof}[Proof of Theorem~\ref{thm:elementary-extension-marked-automorphisms}]
The empty case is immediate, so suppose $D$ is nonempty.
If two graph formulas define the same map over $\mathcal M$, the
sentence asserting their equality remains true over $\mathcal N$.
Functionality, bijectivity, joint continuity in both directions,
parameter preservation, and the marked conditions are first order.
For a setwise mark the condition is equality of its image with that
mark, equivalently membership in the mark if and only if the image
belongs to it.  Thus no closedness assumption on the marks enters.
Composition, graph reversal for inversion, and affine pullback
commute with interpretation.  This proves that the extension map
is a well-defined homomorphism of simplicial groups.

We use the classical finite-simplex description of their homotopy
groups.  For a simplicial group $G$, put
\[
 N_qG=\bigcap_{i=1}^q\ker d_i,\qquad
 Z_qG=\bigcap_{i=0}^q\ker d_i\quad(q\geq1),\qquad Z_0G=G_0.
\]
Then $\pi_q(|G|,1)=Z_qG/d_0N_{q+1}G$, including the component
group for $q=0$.  Simplicial groups are Kan; the multiplication and
normal-boundary assertions needed for this formula are the classical
results in \cite[Theorem~16.6, p.~65; Theorem~17.1 and
Propositions~17.2--17.4, pp.~67--69]{May1967Simplicial}.
Reversing the simplex vertices converts that source's last-face
Moore convention into the displayed first-face convention.  In particular, a positive-degree
class has a single simplex representative equal to the identity
on every face.

\medskip\noindent\textit{Surjectivity.}
Take such a representative $h$ over $\mathcal N$ in degree $q$,
where $q=0$ is allowed.  Regard it as a homeomorphism of the closed
bounded set $E=D\times\Delta^q$.  Choose one graph formula for $h$
and replace its finitely many additional parameters by variables
$a$.  In the coefficient set $T$ impose parameter preservation, all
marked conditions, and, when $q\geq1$, the identity on every face
of $D\times\Delta^q$.  These are finite first-order requirements.
The tuple defining $h$ belongs to $C(N)$ for one cell of the
partition in Lemma~\ref{lem:coefficient-cells-joint-evaluation}.

Every such cell is definably homeomorphic over $\mathcal M$ to
$M^d$ for its dimension $d$
\cite[Proposition~2.5, pp.~18--19]{Coste1999}.
Choose $a_0\in C(M)$ and, in the extended coordinate chart, join
the tuple defining $h$ to $a_0$ by a straight segment.  Pulling back
the joint evaluation and inverse evaluation along this definable
path gives a total homeomorphism on
$D\times\Delta^q\times[0,1]$, preserving all parameters and the
required marks and fixing $D\times\partial\Delta^q\times[0,1]$.
At the last endpoint it is the extension of a simplex over
$\mathcal M$.  Triangulating the parameter prism by finitely many
rational simplices turns this family into a simplicial homotopy;
the inverse restricts on the same simplices.  For $q=0$ it is simply
a path between the two vertices.  Thus the two representatives have
the same homotopy class, proving surjectivity in every degree.

\medskip\noindent\textit{Injectivity.}
Let $h\in Z_qG_{\mathcal M}$ have trivial image in the homotopy
group over $\mathcal N$.  The displayed Moore description gives
$w\in N_{q+1}G_{\mathcal N}$ with $d_0w=h_{\mathcal N}$.
Choose one formula for the graph of $w$ and replace its extra
parameters by a finite variable tuple.  Being a valid marked total
homeomorphism, satisfying $d_iw=1$ for $1\leq i\leq q+1$, and
satisfying $d_0w=h$ are one first-order condition with parameters
in $\mathcal M$.  The condition is satisfiable in $\mathcal N$,
hence in $\mathcal M$.  Its solution is a simplex
$w_0\in N_{q+1}G_{\mathcal M}$ with $d_0w_0=h$.  This proves
injectivity, also for $q=0$.  Translation deals with every other
base component.  The extension map is therefore a weak equivalence.

Apply this result with $D=B\times[0,1]$ and the two marked ends to
obtain the concordance assertion.  For a finite rational PL model,
extending a PL map and extending its definable graph give exactly
the same map.  The resulting square commutes, and its PL vertical
map is a weak equivalence by Theorem~\ref{result-40}.  Its other
vertical map has just been proved to be a weak equivalence.
The assertions about each comparison homomorphism follow.
\end{proof}

\begin{corollary}\label{cor:complete-theory-concordance-comparison}
Fix a complete o-minimal theory in the given language having a real
model, and a rational realization $B$ of a finite abstract closed PL
manifold.  In each homotopy degree, the
injectivity, surjectivity, and bijectivity properties of its
PL-to-definable concordance comparison are the same over all models
of that theory, including the real model.
\end{corollary}

\begin{proof}
Two models of a complete theory have elementary embeddings into a
common model.  Indeed, place their elementary diagrams in disjoint
sets of added constants.  A finite part of either diagram is the
instance of an existential sentence consistent with the complete
theory, hence implied by it.  The two finite parts can therefore be
realized simultaneously.  Compactness gives the common model and
the elementary embeddings; these classical inputs are
\cite[Theorem~2.1.4, pp.~34--35; Lemma~2.3.3, p.~44]{Marker2002ModelTheory}.  The rational PL model is preserved by
both embeddings.  Applying the theorem to the resulting zigzag
proves the assertion.
\end{proof}

The real-model hypothesis is substantive.  In the language of
ordered fields with one named constant, the complete theory of a
real closed field containing a named positive infinite element is
o-minimal but has no real model.  The corollary also does not supply
a comparison property which has not been proved over its real model.

\begin{corollary}\label{cor:countable-definable-automorphism-model}
If the language is countable, the marked automorphism group in
Theorem~\ref{thm:elementary-extension-marked-automorphisms} has a
countable simplicial group as a weak model.
\end{corollary}

\begin{proof}
Start with the finitely many parameters defining $D$ and its marks.
Repeatedly close a countable set under the language's functions and
adjoin a witness to every existential formula with parameters in
the current set which is true in $\mathcal M$.  The union of this
countable sequence is a countable elementary substructure by the
Tarski--Vaught criterion
\cite[Proposition~2.3.5, p.~45]{Marker2002ModelTheory}.  The set and its marks are defined over it,
and the assertion that $D$ is closed and bounded transfers.  There are only
countably many formulas and finite parameter tuples in each
simplicial degree.  The theorem applied to this elementary
substructure proves the claim.
\end{proof}

This last conclusion concerns cardinality.  An arbitrary elementary
diagram need not be decidable, so it supplies neither an effective
test for homeomorphism codes nor an oracle-free recursive
presentation of the homotopy groups.

Countability alone also gives no PL-to-definable weak equivalence.
